\chapter{Contact tracing}
\index{Contact tracing}

\section*{Definition and Overview}
Contact tracing is a fundamental public health intervention used to control the spread of communicable diseases within a population. It is the process of identifying, assessing, and managing individuals (contacts) who have been exposed to an infectious pathogen by a person diagnosed with a contagious disease (the index case). By rapidly locating exposed individuals, notifying them of their risk, and facilitating testing, quarantine, or post-exposure prophylaxis, contact tracing aims to interrupt chains of transmission. This clinical and epidemiological tool is critical for managing endemic infections (such as sexually transmitted infections and tuberculosis), controlling localized outbreaks (such as meningococcal disease or measles), and mitigating the impact of global epidemics and pandemics.

\section*{Epidemiology}
While contact tracing is an intervention rather than a clinical disease, its epidemiological application depends on the dynamics of target pathogens. The efficacy of contact tracing is measured by public health metrics such as the contact identification rate, the time from index symptom onset to contact quarantine, and the secondary attack rate. In many countries, contact tracing is mandated under local public health acts for specific notifiable conditions. The scale of contact tracing operations varies; it remains a continuous, focused activity for pathogens like tuberculosis and sexually transmitted infections (STIs) such as syphilis and HIV, but can be rapidly scaled into a mass community intervention during pandemics, such as influenza or COVID-19 outbreaks, to reduce the effective reproduction number ($R$) of the pathogen.

\section*{Aetiology and Pathophysiology}
The public health decision to initiate contact tracing is based on the biological characteristics of the pathogen and the epidemiological risk of transmission.

\begin{itemize}
    \item \textbf{Transmission dynamics:} Pathogens spread via direct contact, respiratory droplets, airborne aerosols, or sexual contact where proximity or physical intimacy is a prerequisite for transmission.
    \item \textbf{Incubation period:} The specific timeframe between pathogen exposure and the onset of symptoms, during which identifying a contact allows public health officials to intervene before the individual becomes infectious to others.
    \item \textbf{Asymptomatic and presymptomatic shedding:} Pathogens that can replicate and be transmitted by individuals who do not exhibit clinical symptoms, making active contact tracing necessary to uncover silent transmission chains.
    \item \textbf{Pathogen virulence and severity:} Diseases associated with high morbidity, mortality, or long-term sequelae (e.g., Ebola virus disease, measles, invasive meningococcal disease) that justify resource-intensive tracing and containment.
    \item \textbf{Availability of preventive or therapeutic interventions:} Conditions where early detection of contacts allows for the administration of post-exposure vaccination, antibiotics, or antivirals to prevent disease development.
\end{itemize}

\section*{Clinical Features}
Contact tracing is structured into distinct, systematic phases that require clear communication, clinical assessment, and epidemiological tracking.

\begin{itemize}
    \item \textbf{Index case interview:} Conducting a detailed clinical interview with the diagnosed individual to determine their exact symptom onset, infectious window, and the identity and contact information of all close associates.
    \item \textbf{Exposure risk classification:} Categorising contacts based on the duration, proximity, and environmental context of their exposure to the index case (e.g., household contacts vs. casual community contacts).
    \item \textbf{Contact notification:} Reaching out to identified contacts to inform them of their potential exposure while maintaining the strict confidentiality and privacy of the index case.
    \item \textbf{Clinical screening and education:} Assessing contacts for early signs of disease, providing instructions on self-quarantine, and detailing procedures for seeking diagnostic testing safely.
    \item \textbf{Monitoring and follow-up:} Maintaining contact with high-risk individuals throughout the duration of the pathogen's incubation period to track symptom development and ensure adherence to public health directives.
\end{itemize}

\section*{Differential Diagnosis}
In contact tracing, "differential diagnosis" involves the clinical and epidemiological stratification of exposure risk to determine the necessity and intensity of public health interventions.

\begin{itemize}
    \item \textbf{High-risk close contact:} Face-to-face contact or sharing an enclosed space (e.g., a household, classroom, or workplace) with an infectious case for a prolonged period without appropriate personal protective equipment (PPE), requiring immediate quarantine and testing.
    \item \textbf{Low-risk casual contact:} Brief or transient exposure in a well-ventilated or outdoor environment, where the risk of transmission is negligible, requiring symptom monitoring but no restriction of movement.
    \item \textbf{Occupational exposure with protective barriers:} Interactions in healthcare or laboratory settings where standard precautions and appropriate, fitted PPE (such as N95 respirators, gowns, and eye protection) were consistently utilised, classifying the worker as low-risk.
    \item \textbf{Community transmission without known exposure:} Cases arising in individuals with no identifiable contact link, indicating circulating pathogen in the community and the need for broad public testing rather than targeted contact tracing.
    \item \textbf{Non-infectious exposure window:} Contacts whose interaction with the index case occurred outside the calculated infectious period, requiring no public health intervention.
\end{itemize}

\section*{When to Seek Medical Care}
Individuals who are notified that they are a contact of an infectious case should follow the specific public health guidelines provided to them. They should seek medical advice if they develop any symptoms of the disease, notifying the medical facility of their contact status prior to arrival.

\textbf{Contact your local emergency services for:}
\begin{itemize}
    \item Severe shortness of breath, chest pain, or bluish discolouration of the lips (indicating respiratory failure from pathogens like COVID-19 or influenza).
    \item Sudden onset of high fever accompanied by a stiff neck, severe headache, confusion, and a non-blanching dark red or purple rash (suggestive of meningococcal disease).
    \item Severe lethargy, confusion, or loss of consciousness in an exposed individual.
    \item Rapidly worsening systemic illness or signs of sepsis.
\end{itemize}

\section*{Diagnosis and Investigations}
Diagnostic testing is directed at confirming infection in both the index case and the identified contacts to guide clinical management and tracing expansion.

\begin{itemize}
    \item \textbf{Nucleic acid amplification tests (NAAT):} Polymerase chain reaction (PCR) assays used to detect pathogen-specific DNA or RNA, representing the gold standard for diagnosing respiratory (e.g., influenza, pertussis) and sexually transmitted infections.
    \item \textbf{Serology and antibody testing:} Blood tests to detect IgM and IgG antibodies, useful for determining past exposure, current infection stage, or immunity status (e.g., for syphilis, HIV, or hepatitis).
    \item \textbf{Microbiological culture:} Isolating viable bacteria from swabs, blood, or sputum specimens to confirm infection and perform antimicrobial susceptibility testing.
    \item \textbf{Antigen detection assays:} Rapid diagnostic tests (such as rapid antigen tests) that detect specific viral proteins, offering rapid, point-of-care results during outbreaks.
    \item \textbf{Radiological imaging:} Chest X-rays or CT scans performed on contacts of active pulmonary tuberculosis to screen for subclinical lung lesions.
\end{itemize}

\section*{Management}
The management of contacts aims to prevent secondary transmission, monitor for disease onset, and provide prophylaxis when clinically indicated.

\begin{itemize}
    \item \textbf{Quarantine and isolation:} Restricting the movement of asymptomatic exposed individuals for the duration of the incubation period to prevent community spread in case they develop the disease.
    \item \textbf{Post-exposure prophylaxis (PEP):} Administering antimicrobial agents or vaccines immediately after exposure to prevent infection (e.g., ciprofloxacin for meningococcal contacts, post-exposure vaccination for rabies or measles, and antiretrovirals for HIV exposures).
    \item \textbf{Partner notification support:} Providing resources, such as patient-delivered partner therapy (PDPT) or anonymous web-based notification tools, to assist individuals diagnosed with STIs in notifying their sexual partners.
    \item \textbf{Digital contact tracing tools:} Utilizing electronic databases, contact tracing software, and mobile applications to track exposures, map transmission networks, and automate notification alerts.
    \item \textbf{Symptomatic and supportive care:} Advising contacts on self-care, hydration, rest, and appropriate antipyretics while they remain in quarantine.
\end{itemize}

\section*{Complications}
Challenges and complications arising from contact tracing can impact both public health containment and individual wellbeing.

\begin{itemize}
    \item \textbf{Transmission leakage:} Delays or failure to identify and isolate high-risk contacts, leading to unchecked community transmission and outbreak expansion.
    \item \textbf{Psychosocial impact of quarantine:} Increased rates of anxiety, depression, loneliness, and financial stress experienced by individuals undergoing mandatory isolation.
    \item \textbf{Confidentiality and privacy breaches:} Accidental disclosure of an index case's identity or health status, leading to social stigma, discrimination, or loss of trust in healthcare providers.
    \item \textbf{Public health workforce overload:} High case volumes exceeding the capacity of public health departments, resulting in delayed interviews, backlogs, and ineffective containment.
    \item \textbf{Public fatigue and non-compliance:} Resistance to contact tracing due to pandemic fatigue, distrust of government, misinformation, or fear of punitive isolation measures.
\end{itemize}

\section*{Prognosis and Outcomes}
The prognosis of a contact tracing initiative is determined by its timeliness, coverage, and the cooperation of the affected community. When contact tracing is implemented rapidly and achieves high coverage (e.g., tracing and quarantining >80\% of contacts within 48 hours), it can successfully suppress outbreaks and prevent endemic establishment. For individual contacts, the prognosis is significantly improved by tracing, as it facilitates early diagnosis, close clinical monitoring, and the prompt administration of treatments or prophylactic therapies, reducing the risk of severe clinical outcomes or complications.

\section*{Prevention and Self-Care}
Individuals can support contact tracing and protect themselves by adhering to public health recommendations and practicing proactive health measures.

\begin{itemize}
    \item \textbf{Maintaining awareness of exposures:} Keeping track of personal close contacts and noting environments visited during active community outbreaks.
    \item \textbf{Prompt testing and reporting:} Seeking diagnostic testing immediately upon developing symptoms or when notified of a potential exposure, and reporting positive results to health authorities.
    \item \textbf{Strict compliance with isolation:} Remaining at home and avoiding physical contact with household members and the community during mandated quarantine periods.
    \item \textbf{Honest communication with health workers:} Cooperating fully with contact tracers by providing accurate details regarding close contacts and exposure timelines.
    \item \textbf{Adoption of preventative behaviours:} Wearing masks in crowded indoor spaces, practicing hand hygiene, and maintaining up-to-date immunisations to minimise transmission risk.
\end{itemize}

\section*{Sources and Further Reading}
The following reputable organisations provide further information and clinical guidelines on contact tracing:

\begin{itemize}
    \item World Health Organization. \textit{Operational guidance on contact tracing for infectious diseases and outbreak management.} https://www.who.int/
    \item Australian Government Department of Health and Aged Care. \textit{National guidelines for public health units on contact tracing and notifiable diseases.} https://www.health.gov.au/
    \item Healthdirect Australia. \textit{Public advice on contact tracing, isolation requirements, and testing locations.} https://www.healthdirect.gov.au/
    \item Centers for Disease Control and Prevention. \textit{Resources, training materials, and digital tools for contact tracing workflows.} https://www.cdc.gov/
    \item European Centre for Disease Prevention and Control. \textit{Epidemiological resources and frameworks for contact tracing in communicable disease control.} https://www.ecdc.europa.eu/
\end{itemize}
